In this paper, we introduce the Tamarin prover following \cite{TamarinThesisInduction}.
We omit certain details for clarity.
For full details, we refer the reader to \cite{TamarinThesisInduction}.
In Section~\ref{sec:preliminaries}, we introduce preliminaries that are independent of Tamarin itself: order-sorted term algebras, term rewriting, equational theories, and multisets.
In Sections~\ref{subsec:protocol-specs}, \ref{subsec:semantics}, and~\ref{sec:ndgs} we introduce Tamarin itself.
We start with Tamarin's syntax, then introduce Tamarin's semantics, and finally introduce how variants determine Tamarin's reasoning.

Tamarin works in the \emph{symbolic model}, where protocol messages are represented as elements of a term algebra $\Sigma$.
An equational theory encodes the semantics of these terms.
Along with the equational theory, protocol and attacker behavior are specified by a set of multiset rewriting rules.
A Tamarin protocol model is a pair of an equational theory and a set of multiset rewriting rules.
The semantics of a Tamarin model are its set of associated traces, i.e., sequences of multiset rewriting rule applications.
For our optimization, we will require the results of Theorem~\ref{thm:reasoning-mod-ax}, which shows that reasoning about a set of protocol model rules modulo a full equational theory is equivalent to reasoning about a set of protocol model rule \emph{variants} modulo associativity-commutativity (if applicable).

\subsection{Preliminaries}
\label{sec:preliminaries}

\subsubsection{Order-Sorted Term Algebras}

An \emph{order-sorted signature} $\Sigma = (S, \leq, F)$ consists of a set of sorts $S$, a partial order $(\leq) \subseteq S \times S$ on the sorts, and a set of function symbols $F$.
We write $f/n$ for a function symbol $f \in F$ with arity $n$.
Every sort $s \in S$ has an infinite set of associated variables $\vars_s$ and constants $\consts_s$, which are all pairwise disjoint.
We write $\ofsort{v}{s}$ for $v \in \vars_s$ and $\vars = \bigcup_{s \in S} \vars_s$ for the set of all variables.
Functions in $F$ have the form $f: s_1 \times \dots \times s_n \rightarrow s$ (where $s_1, \dots, s_n, s \in S$).
For $s \in S$, we use $\terms_\Sigma^s$ to denote the terms of sort $s$.
The set of all terms of an order-sorted signature is $\terms_\Sigma = \bigcup_{s \in S} \terms_\Sigma^s$.
A term is \emph{ground} if it contains no variables.

A substitution is a function $\sigma: \vars \to \terms_\Sigma$ mapping variables to terms.
We require that substitutions respect sorts, i.e., variables of sort $s$ must be mapped to terms of sort $s' \leq s$.
We extend substitutions homomorphically to terms and other constructs (to be defined later) as usual, and write $t\sigma$ for an application of $\sigma$ to a term $t \in \terms_\Sigma$.
A \emph{valuation} is a substitution to ground terms.
Valuations can also be defined as explicit lists.
In that case, variables not mapped in that list are irrelevant and can map to arbitrary terms.

\subsubsection{Term Rewriting}

A \emph{position} $p$ in a term is a list of indices.
We write $t|_p$ for \emph{$t$'s subterm at position $p$}, and define it inductively as:
\begin{equation*}
  t|_p := \begin{cases}
    t & p = [] \\
    t_{i_0}|_{[i_1, \dots, i_n]} & p = [i_0, i_1, \dots, i_n], t = f(t_1, \dots, t_k), 1 \leq i_0 \leq k \\
    \bot & \text{otherwise}
  \end{cases}
\end{equation*}

We write $t[s]_p$ to denote the term obtained from replacing $t_p$ in $t$ with $s$.
We say that $t$ is a \emph{syntactic subterm} of $t'$ if there exists a position $p$ such that $t'|_p = t$.
Whenever this is the case, we write $t \sqsubseteq t'$.

A \emph{rewrite rule} is an ordered pair of terms $l \to r$.
A \emph{rewriting system} $\R$ is a set of rewrite rules.
The corresponding \emph{rewrite relation} $\to_R$ is defined as $s \to_R t$ if there is a position $p$ in $s$, a rewrite rule $l \to r \in \R$, and a substitution $\sigma$ such that $s|_p = l\sigma$ and $t = s[r\sigma]_p$.

A rewriting system $\R$ is \emph{terminating} if there is no infinite sequence of terms $(t_i)_{i \in \nat}$ such that $t_i \to_R t_{i+1}$ ($i \in \nat$).
$\R$ is \emph{confluent} when for all terms $t$, $s_1$, and $s_2$, if $t \to_R^* s_1$ and $t \to_R^* s_2$ then there exists a term $t'$ such that $s_1 \to_R^* t'$ and $s_2 \to_R^* t'$.
$\R$ is \emph{convergent} if it is terminating and confluent.
We use $t{\downarrow}_R$ to denote a term $t$'s normal form in a convergent rewriting system $\R$, which is the term $t'$ such that $t \to_R^* t'$ and there exists no term $t''$ such that $t' \to_R t''$.
Observe that every term's normal form is uniquely determined.

\subsubsection{Equational Theories}

Protocol models are defined with respect to an order-sorted signature $\Sigma$ and a set of equations $E$ that have the form $s \simeq t$ for $s, t \in \terms_\Sigma$.
Typically, the equations model the relationships between different cryptographic operations, e.g., that decryption is the inverse operation to encryption.
$\Xi = (\Sigma, E)$ is an \emph{equational presentation}.
An \emph{equational theory} $=_\Xi$ is the smallest congruence relation over $\Sigma$ containing all instances of $E$.
We often leave $\Sigma$ implicit and identify $\Xi$ by $E$.
Two terms $s$ and $t$ are \emph{equal modulo $E$} if $t =_E s$.
We will use the $E$-subscript to denote operators modulo $E$, e.g., $\in_E$.

If an equational theory $E$'s equations can be oriented to a convergent rewriting system $\R$, then $t =_E s$ if and only if $t{\downarrow}_R = s{\downarrow}_R$.
Some equations, e.g., commutativity, cannot be oriented and in those cases, we split an equational theory $E$ into oriented rewriting rules $\R$ and unoriented equations $AX$.
The \emph{rewrite relation} $\to_{\R,AX}$ is defined as $s \to_{\R,AX} t$ if there is a position $p$, a rewrite rule $l \to r \in \R$, and a substitution $\sigma$ such that $s|_p =_{AX} l\sigma$ and $t = s[r\sigma]_p$.
We say that $\R$ is $AX$-convergent if $\to_{\R,AX}$ is terminating and confluent.
We write $t\normf^\R$ to denote $t$'s normal form with respect to $\to_{\R,AX}$.
$\R$ is $AX$-coherent when for all terms $t_1$, $t_2$, and $t_3$, if $t_1 \to_{\R,AX}^* t_2$ and $t_1 =_{AX} t_3$ then there exist terms $t_4$, $t_5$ such that $t_4 =_{AX} t_5$ and $t_2 \to^*_{\R,AX} t_4$ and $t_3 \to_{\R,AX}^+ t_5$.
We write $\R^\simeq := \{ l \simeq r \mid l \to r \in \R \}$ for a rewriting system $\R$'s associated equational theory.
For $E = \R^\simeq \cup AX$, if $\R$ is $AX$-convergent and $AX$-coherent, then $t =_E s$ if and only if $t\normf^\R =_{AX} s\normf^\R$.

\subsubsection{Multisets}

A \emph{multiset} $\{a, b, \dots\}^\sharp$ is a set of items $a$, $b$, etc., where each item can occur multiple times.
We extend the usual operations and constants for sets to multisets using the $\sharp$ superscript, e.g., $\emptyset^\sharp$ is the empty multiset.
Whenever we use relations, e.g., $\subseteq^\sharp$, or functions, e.g., $\cup^\sharp$, on multisets, they take cardinality into account.

\subsection{Protocol Specifications}
\label{subsec:protocol-specs}

In Tamarin, state transitions are modeled using multiset rewriting rules.
The global state consists of a multiset of ground facts, and the rewriting rules define how the state is updated.
Facts are similar to predicates in first-order logic and are of the form $F(t_1, \dots, t_n)$.
$F$ is a fact symbol from a fact signature $\factsig$, and $t_1, \dots, t_n$ are terms defined over the model's signature $\Sigma$.
We fix the sorts in $\Sigma$ to $\msg$, $\fresh$, $\pub$, and $\sortnat$.
$\fresh$, $\pub$, and $\sortnat$ are subsorts of $\msg$.
The sort $\fresh$ models fresh, unguessable values such as cryptographic random numbers.

A fact is ground if all its terms are ground.
Facts typically model the local state of participants or the network, e.g., who possesses which keys, or what messages have been sent over the network.
We consider the special, reserved fact symbols $\Fr$, $\In$, $\Out$, $\Kup$, and $\Kdown$ of arity 1 in $\factsig$.
These symbols respectively model generating a fresh, unguessable value, receiving and sending a message over an insecure network, and adversary reasoning.

Multiset rewriting rules in Tamarin have the form $$[l_1, \dots, l_m] \lrule{a_1, \dots, a_n} [r_1, \dots, r_o].$$
$l_i$, $a_j$, $r_k$ are facts and may use free variables ($1 \leq i \leq m$, $1 \leq j \leq n$, $1 \leq k \leq o$).
$l_1, \dots, l_m$ are the rule's \emph{premises}, $a_1, \dots, a_n$ the rule's \emph{actions} (also called \emph{action facts}), and $r_1, \dots, r_o$ the rule's \emph{conclusions}.
$[a_1, \dots, a_n]$ can be omitted when a rule has no actions.
Rules can be instantiated by a substitution $\sigma$.
One rule $ri$ is another rule $r$'s \emph{instance} if there exists a substitution $\sigma$ such that $r\sigma = ri$.
We write $\prems(r)$, $\acts(r)$, and $\concs(r)$ for a rule $r$'s list of premises, actions, and conclusions respectively.
A rule (instance) is ground if all its facts are ground.

Facts are either persistent (preceded by a $!$) or linear.
When a rule is applied, the linear facts in its premise are removed from the global state and replaced by the facts in its conclusion.
Persistent facts are not consumed when used in a rule's premise.

A set of multiset rewriting rules $P$ is a \emph{multiset rewriting system} if
\begin{enumerate}[label=(\roman*)]
  \item no rule uses a fresh constant,
  \item no rule's premise contains an $\Out$, $\Kup$, or $\Kdown$ fact, and
  \item no rule's conclusion contains a $\Fr$, $\In$, $\Kup$, or $\Kdown$ fact, and
  \item all $\fresh$, $\msg$, and $\sortnat$ variables used in a rule's conclusion are introduced in a rule's premise.
\end{enumerate}

\paragraph{Built-In Rules}

Tamarin uses a number of built-in rules to model the generation of random values and the adversary deduction following the model's equational theory $E$, the latter of which we refer to by $ND$ (normal form message deduction rules).
For this work, the following two built-in rules are particularly relevant:
\begin{gather*}
  \RFresh := [] \rightarrow [ \Fr(\ofsort{x}{\fresh}) ],~\text{and} \\
  [ \Fr(\ofsort{x}{\fresh}) ] \lrule{\Kup(\ofsort{x}{\fresh})} [ \Kup(\ofsort{x}{\fresh}) ] \in ND.
\end{gather*}

For a full introduction of all built-in rules see \cite{TamarinThesisInduction}.

\subsection{Semantics}
\label{subsec:semantics}

\subsubsection{Executions and Traces}

Given a multiset rewriting system $P$, adversary deduction rules $ND$, an equational theory $E$, and $R = P \cup ND \cup \{ \RFresh \}$, the tuple $(R,E)$ is a \emph{Tamarin} or \emph{protocol model}.
Tamarin models induce a labeled state transition system.
The initial state is the empty multiset $\emptyset^\sharp$.
A state transition is the application of a ground rule instance $ri = l \lrule{a} r$.
We use $\lfacts(x)$ and $\pfacts(x)$ to denote the linear and persistent facts respectively in $x$.
$ri$ can be applied if its premises are contained in the global state, which is then updated accordingly.
Formally, the state transition relation $\Rightarrow_{(R, E)}$ is defined as:
\begin{align*}
  S \Rightarrow_{(R, E)} (S \setminus^\sharp \lfacts(l)) \cup^\sharp r && \text{if}~\lfacts(l) \subseteq^\sharp S, \pfacts(r) \subseteq \setof(S)
\end{align*}

An \emph{execution} is a sequence of labeled state transitions.
We only consider executions where each instance of $\RFresh$ is unique.
% TODO: Intuition for what this model
Each execution has a corresponding \emph{trace}, which is the list of sets of action facts labeling each transition.
The semantics of a Tamarin model is its set of traces.

\subsubsection{Dependency Graphs}

Dependency graphs capture the sequence of rule applications and dependencies between rule premises and conclusions.
They are closely related to the constraint systems used in Tamarin's protocol analysis, and they enable effective proof search.

\begin{definition}[Dependency graph] \label{def:dependency-graph}
  Let $(R, E)$ be a Tamarin model.
  A \emph{dependency graph} is a tuple $dg = (I, D)$ where $I$ is a finite list of ground rule instances in $R$, and $D \subseteq \nat^2 \times \nat^2$.
  A rule's index in $I$ is its \emph{timepoint}.
  We write $(i,u) \edge (j,v)$ for $((i,u), (j,v)) \in D$ when $D$ is clear from the context.
  Here, $(i, u)$ denotes the conclusion with index $u$ of the rule with index $i$ in $I$.
  Similarly, $(j,v)$ denotes the premise with index $v$ of the rule with index $j$.
  We require that edges correctly connect premises and conclusions, i.e., $i < j$ and the conclusion $(i, u)$ is equal modulo $E$ to the premise $(j, v)$.
  Furthermore, each premise must have exactly one incoming edge, every linear conclusion has at most one outgoing edge, and instances of $\RFresh$ are unique.
\end{definition}

\begin{definition}[Trace]
  A \emph{trace} induced by a dependency graph $dg = (I, D)$ is a list $\trace(dg)$ of sets of action labels.
  Concretely, $\trace(dg)$ is of length $|I|$ and every index $i$ corresponds to the set of action labels at timepoint $i$, i.e., to $\setof(\acts(I_i))$.
\end{definition}

One can show that the dependency graphs of a model $(R,E)$ induce the same set of traces as the set of traces derived from the model's executions (\cite[Theorem 3]{TamarinThesisInduction}).
We write $\dgraphs_E(R)$ for all dependency graphs of a protocol model $(R, E)$ and $\trace(dg)$ for a dependency graph's associated trace.

\subsubsection{Trace Properties}

In Tamarin, protocol properties are expressed as first-order logic \emph{trace properties}, which are given by closed \emph{trace formulas}.
In the following, we briefly relate trace properties to protocol models, but do not fully define them and their semantics.
We refer to \cite{TamarinThesisInduction} for a complete definition.

Given a valuation $\theta$ of the free variables of $\varphi$, we write $(tr, \theta) \tsatE \varphi$ to mean that the trace $tr$ satisfies $\varphi$ under the valuation $\theta$.
A trace formula $\varphi$ is \emph{valid} for a protocol model $(R, E)$, written $R \tsatE^\forall \varphi$, if for all traces $tr$ of $(R, E)$ and every valuation $\theta$, one has $(tr, \theta) \tsatE \varphi$.
$\varphi$ is \emph{satisfiable} in $(R, E)$, written $R \tsatE^\exists \varphi$, if there exists a trace $tr$ of $(R, E)$ and valuation $\theta$ such that $(tr, \theta) \tsatE \varphi$.

\subsection{Normal Form Dependency Graphs}
\label{sec:ndgs}

Let $E$ be an equational theory, $\R$ be a rewriting system, and $AX \subseteq E$ the equations for associativity and commutativity in $E$ such that
\begin{enumerate*}[label=(\roman*)]
  \item $E = \R^\simeq \cup AX$ and
  \item $\R$ is $AX$-convergent and $AX$-coherent.
\end{enumerate*}
$E$ has the \emph{finite variant property} if for all terms $t$ there is a finite set of substitutions $T = \{\tau_1, \dots, \tau_k\}$ such that for all substitutions $\sigma$ there is a substitution $\tau \in T$ and substitution $\sigma'$ so that $t\sigma \normf^\R =_{AX} (t\tau \normf^\R) \sigma'$ and for all variables $x$ occurring in $t$, we have $x\sigma \normf^\R =_{AX} x\tau\sigma'$.

The set $\variants{t}^\R := \{ (t\tau_i\normf^\R, \tau_i) \mid 1 \leq i \leq k \}$ is a \emph{complete set of $\R,AX$-variants of $t$}.
We extend the notion of complete variants to multiset rewriting rules and multiset rewriting systems straightforwardly.
Intuitively, the variants of a term or rule allow us to reason modulo $E$ by reasoning modulo normal forms and $AX$ only.

\begin{theorem}[Changing the Equational Theory {\cite[p.114]{TamarinThesisInduction}}] \label{thm:reasoning-mod-ax}
  For every guarded\footnote{%
    Tamarin only considers trace properties that are \emph{guarded}, i.e., a subset of all possible trace properties.
    The details of guardedness are irrelevant to this work, and we thus omit a formal definition.
    For a definition, see~\cite{TamarinThesisInduction}.
  } trace property $\varphi$ and protocol model $(R, E)$ where $\R^\simeq \cup AX = E$:
  $$R \tsatE^\forall \varphi \iff \{ \trace(dg) \mid dg \in dgraphs_{AX}(\variants{R}^\R), dg~\text{is}~\normf^\R\text{-normal} \} \tsat{AX}^\forall \varphi.$$
\end{theorem}
